\section{Accesibilidad y microcopy}

\subsection{Descripcion}

Este documento define lineamientos funcionales de accesibilidad y de microcopy para Quiniela. Su objetivo es asegurar que el producto comunique acciones, restricciones, errores y estados con lenguaje claro, accesible y consistente con el comportamiento real del sistema.

\subsection{Alcance}

Aplica a textos de botones, mensajes de estado, validaciones, errores, confirmaciones y accesibilidad basica de la interfaz. No define una guia editorial completa de marca.

\subsection{Definiciones}

\begin{itemize}
  \item \textbf{Microcopy}: texto corto orientado a guiar acciones, estados o errores.
  \item \textbf{Accesibilidad basica}: minimo de buenas practicas para que la interfaz sea comprensible y operable por una audiencia amplia.
\end{itemize}

\subsection{Reglas}

\begin{itemize}
  \item El lenguaje debe ser simple, directo y no tecnico para el participante.
  \item Los mensajes del panel administrativo pueden ser mas precisos, pero no ambiguos ni cripticos.
  \item La accesibilidad no puede depender solo del color.
  \item Ningun mensaje debe ocultar una regla real del sistema o sugerir una accion inexistente.
\end{itemize}

\subsubsection{Principios de microcopy}

\begin{itemize}
  \item describir primero el estado o problema y luego la accion siguiente;
  \item evitar tecnicismos cuando el usuario no los necesita;
  \item usar la misma palabra para el mismo concepto en todas las pantallas;
  \item no prometer exito hasta que la operacion haya sido confirmada.
\end{itemize}

\subsubsection{Mensajes funcionales prioritarios}

\begin{itemize}
  \item cuenta pendiente de activacion;
  \item comprobante recibido;
  \item activacion aprobada o rechazada;
  \item pick guardado correctamente;
  \item pick bloqueado por cierre;
  \item partido sin pick al llegar el lock;
  \item prediccion especial bloqueada por cierre;
  \item resultado corregido y puntajes recalculados;
  \item accion administrativa realizada o rechazada.
\end{itemize}

\subsubsection{Lineamientos de accesibilidad minima}

\begin{itemize}
  \item los estados importantes deben apoyarse en texto y no solo en color o icono;
  \item los CTA deben tener etiquetas claras y especificas;
  \item los errores de formulario deben aparecer cerca del campo o accion afectada;
  \item la jerarquia visual debe permitir identificar rapidamente titulo, accion principal y estado relevante;
  \item las tablas o rankings deben seguir siendo interpretables cuando cambie el orden o exista recalculo.
\end{itemize}

\subsubsection{Tono recomendado}

\begin{itemize}
  \item claro;
  \item respetuoso;
  \item orientado a accion;
  \item transparente cuando el sistema corrige o restringe algo.
\end{itemize}

\subsubsection{Mensajes que deben evitarse}

\begin{itemize}
  \item mensajes vagos como ``ocurrio un problema'' sin contexto util;
  \item mensajes que culpen al usuario cuando el fallo es tecnico;
  \item textos que sugieran que un pick fue guardado si la operacion aun no se confirmo;
  \item mensajes administrativos que oculten impacto real de una accion critica.
\end{itemize}

\subsection{Edge Cases}

\begin{itemize}
  \item Un error tecnico puede requerir lenguaje diferente al error funcional aunque aparezca en la misma pantalla.
  \item El aviso de resultado corregido debe ser claro, pero no necesita exponer por defecto detalle profundo de auditoria.
  \item Un estado de cuenta bloqueada requiere tono firme y claro sin dejar ambiguedad sobre la restriccion.
\end{itemize}

\subsection{Invariantes}

\begin{itemize}
  \item Los mensajes no ocultan reglas reales del sistema.
  \item La accesibilidad no depende solo del color.
  \item El microcopy mantiene consistencia terminologica con la documentacion del dominio.
\end{itemize}

\subsection{Preguntas abiertas}

\begin{itemize}
  \item Confirmar si el proyecto quiere un glosario corto de microcopy prohibido y recomendado para mantener consistencia entre frontend y operacion.
\end{itemize}

\subsection{Decisiones pendientes}

\begin{itemize}
  \item \texttt{Decision pendiente} \texttt{No bloqueante}: definir un catalogo operativo de mensajes clave por feature si el proyecto quiere controlar microcopy con mayor detalle.
  \item \texttt{Decision pendiente} \texttt{No bloqueante}: confirmar el tono institucional final para mensajes administrativos y comunicaciones visibles al participante.
\end{itemize}
